\chapter{Hardware Implementation, Timing, and Synthesis Evaluation}
\label{chap:synthesis_evaluation}

\section{Implementation Environment and Target Platform}
The Multi-Head Attention hardware accelerator was synthesized, implemented, and placed-and-routed using **AMD Vivado Design Suite 2026.1** targeting the **AMD Kria KV260 Vision AI Starter Kit** (\texttt{xck26-sfvc784-2LV-c}). 

\begin{itemize}
    \item \textbf{Target Operating Frequency:} 200.0~MHz ($T_{\text{clk}} = 5.000$~ns).
    \item \textbf{Synthesis Mode:} Out-of-Context (OOC) and Full Block Design integration with AXI DMA and Zynq UltraScale+ MPSoC.
    \item \textbf{Constraints:} Defined in \texttt{timing\_kv260.xdc} with 0.500~ns input/output port budgets.
\end{itemize}

\section{Post-Routing Hardware Resource Utilization}
Table~\ref{tab:post_route_util} summarizes the post-routing resource consumption extracted directly from Vivado implementation reports (\texttt{post\_route\_utilization.rpt}).

\begin{table}[htbp]
    \centering
    \caption{Post-Routing Resource Utilization on AMD Kria KV260 (Variant 2A)}
    \label{tab:post_route_util}
    \begin{tabularx}{\textwidth}{l r r r}
        \toprule
        \textbf{Resource Primitive} & \textbf{Used} & \textbf{Available (XCK26)} & \textbf{Utilization (\%)} \\
        \midrule
        \textbf{CLB LUTs}               & \textbf{10,631} & 117,120 & \textbf{9.08\%} \\
        • \textit{LUT as Logic}         & 5,891           & 117,120 & 5.03\% \\
        • \textit{LUT as Memory (RAM)}  & 4,736           & 57,600  & 8.23\% \\
        • \textit{LUT as Shift Register}& 4               & 57,600  & 0.01\% \\
        \textbf{CLB Registers (FF)}     & \textbf{7,366}  & 234,240 & \textbf{3.14\%} \\
        \textbf{DSP48E2 Slices}         & \textbf{52}     & 1,248   & \textbf{4.17\%} \\
        • \textit{Stage 1 QK$^T$ GEMM}  & 8               & 1,248   & 0.64\% \\
        • \textit{Stage 2 Dyadic Scaler}& 1               & 1,248   & 0.08\% \\
        • \textit{Stage 3/4 \& Top Core}& 43              & 1,248   & 3.45\% \\
        \textbf{Block RAM (RAMB36E2)}   & \textbf{1}      & 144     & \textbf{0.69\%} \\
        \textbf{UltraRAM (URAM)}        & \textbf{0}      & 64      & \textbf{0.00\%} \\
        \textbf{Clock Buffers (BUFGCE)} & \textbf{1}      & 112     & \textbf{0.89\%} \\
        \bottomrule
    \end{tabularx}
\end{table}

\section{Timing Closure Analysis at 200 MHz}
Timing analysis across all process corners confirmed complete timing closure at the full target operating frequency of 200~MHz. The post-routing timing summary (\texttt{post\_route\_timing.rpt}) reports:

\begin{table}[htbp]
    \centering
    \caption{Post-Routing Timing Summary at 200.0 MHz ($T_{\text{clk}} = 5.000$~ns)}
    \label{tab:timing_summary}
    \begin{tabular}{l c c c}
        \toprule
        \textbf{Timing Parameter} & \textbf{Slack Value} & \textbf{Failing Endpoints} & \textbf{Total Endpoints} \\
        \midrule
        Worst Negative Slack (WNS) & \textbf{$+0.051$~ns} & 0 & 51,585 \\
        Total Negative Slack (TNS) & \textbf{$0.000$~ns}  & 0 & 51,585 \\
        Worst Hold Slack (WHS)     & \textbf{$+0.016$~ns} & 0 & 51,585 \\
        Total Hold Slack (THS)     & \textbf{$0.000$~ns}  & 0 & 51,585 \\
        Worst Pulse Width Slack (WPWS) & \textbf{$+1.738$~ns} & 0 & 12,109 \\
        \bottomrule
    \end{tabular}
\end{table}

All 51,585 synchronous endpoints meet setup and hold requirements with zero timing violations, confirming stable hardware operation.

\section{Latency, Throughput, and Real-Time Frame Rate}
Table~\ref{tab:latency_breakdown} details the cycle counts and wall-clock execution time for a full Multi-Head Attention layer ($N = 256, d_k = 16, H = 4$).

\begin{table}[htbp]
    \centering
    \caption{Latency and Throughput Breakdown per Attention Head at 200 MHz}
    \label{tab:latency_breakdown}
    \begin{tabularx}{\textwidth}{l r r X}
        \toprule
        \textbf{Operation} & \textbf{Cycles / Head} & \textbf{Time (@ 200 MHz)} & \textbf{Description} \\
        \midrule
        Ingress DMA Load ($Q, K, V$) & 768   & $3.84\ \mu\text{s}$  & Overlapped via Ping-Pong \\
        Stage 1: $QK^T$ GEMM         & 66,816 & $334.08\ \mu\text{s}$ & $II = 1$ score/cycle \\
        Stage 2: Dyadic Scaler       & —      & —                    & Pipelined inside Stage 1 \\
        Stage 3: IntSoftmax (Pass 1..3)& 70,144& $350.72\ \mu\text{s}$ & Max-Search + Exp + Div \\
        Stage 4: Context GEMM ($P \times V$) & 72,192 & $360.96\ \mu\text{s}$ & 16-channel accumulation \\
        Egress DMA Writeback ($O$)   & 256    & $1.28\ \mu\text{s}$  & Pipelined with Stage 4 \\
        \midrule
        \textbf{Total per Head}      & \textbf{209,919} & \textbf{1.050~ms} & Single head latency \\
        \textbf{Total for Full Layer ($H=4$)} & \textbf{839,676} & \textbf{4.198~ms} & Complete MHA computation \\
        \midrule
        \textbf{Attention Frame Rate}& \multicolumn{3}{c}{\textbf{238.2 Frames Per Second (FPS)}} \\
        \bottomrule
    \end{tabularx}
\end{table}

At 238.2 FPS, the accelerator outperforms the 30 FPS real-time threshold by 7.94$\times$, leaving substantial timing margin for pre- and post-processing networks.

